\begin{rubric}{\textsc{Awards and Grants}}

%% VERIFY & UNCOMMENT once the MINOS outcome is known. Grants are the single
%% highest-value line on an early-career faculty CV, so it is worth listing
%% pending applications explicitly if the decision is still outstanding.
%\entry*[2026]%
%MINOS: `Motives for cooperation across societies' (\texteuro 15,000, 24 months, funding body: [FUNDER]) --- Principal Investigator

\entry*[2022]%
Project Grant for `Crowd-sourcing the quality of wine' (\$2,000, funding body: American Association of Wine Economists)

%\entry*[2018]%
% Project Grant for `Is there an `I' in persuasion' (\pounds3,150, funding body: Economic and Social Research Council (ES/P008976/1))

\entry*[2014 -- 2017]%
	Vice Chancellor's Scholarship for Academic Excellence, University of Nottingham (\pounds43,740)

\entry*[2017]	Economics Teaching Excellence Scholarship, University of Nottingham (\pounds9,000)
			\par Based on 2 consecutive Teaching Excellence Awards (2015, 2016)	

\entry*[2017] Project Grant for `A horse race between elicitation methods of Cumulative Prospect Theory'
(\pounds2,500, funding body: ESRC-funded Network for Integrated Behavioural Science (ES/K002201/1)) 

\entry*[2016] Project Grant for `The Description--Experience Gap in Cooperation' 
(\pounds3,000, funding body: ESRC-funded NIBS (ES/K002201/1))

\entry*[2016] Project Grant for `Limits of the social-benefit motive among high-risk patients: a field 
experiment on influenza vaccination behaviour' (\pounds660, funding body: University of Health Sciences, Istanbul Training and Research Hospital)

\entry*[2015] Interdisciplinary Collaboration Project Award 
(\pounds1,500 project grant for `Seeding a Seed': an interdisciplinary research project on antimicrobial resistance, University of Nottingham)

\entry*[2014]	`Datathlon, Greece 2014' (3rd place): national big-data analysis competition.

\end{rubric}
